\section{Validación integral del esquema}
El archivo DDL.sql se ensambla por fases: creación de todas las tablas,
dominios, claves primarias y candidatas, y finalmente referencias. Dentro de
cada fase se incorporan los módulos 01 a 05. El resultado es SQL autónomo,
sin instrucciones de inclusión ni dependencias de archivos locales.

Se ejecutó en una base nueva con PostgreSQL 18, mediante psql con
\texttt{ON\_ERROR\_STOP=1} y transacción única. La ejecución terminó sin errores.
Las 30 tablas y los nombres literales de sus 146 columnas se contrastaron
con el archivo Draw.io relacional. El catálogo confirma 30 claves primarias,
38 claves foráneas, seis restricciones UNIQUE, 134 CHECK y 146 columnas
NOT NULL. Ninguna restricción quedó sin validar.

Se conservan los rótulos \texttt{Genrente} y \texttt{TeléfonoTécnivo}, que
contienen erratas en el modelo original, para mantener correspondencia literal.
Las tablas de contacto conservan sus propios nombres, por ejemplo
\texttt{CorreoElectrónicoGerente}. La clave de Mantenimiento es compuesta.
En Disponer y Administrar, cuyos atributos aparecen marcados como FK sin
marcas PK explícitas, se define una clave primaria con la pareja de columnas.
Administrar añade UNIQUE por gerente y por sucursal para la relación uno a uno.

La prueba conjunta insertó datos válidos en las 30 tablas reales y comprobó
363 operaciones que debían rechazarse: 146 por NULL, 138 por dominio,
36 por unicidad y 43 por integridad referencial. Incluyó referencias entre
los cinco módulos, pares de tarjeta inexistentes, puestos incorrectos en canjes,
eliminaciones de filas referenciadas y un discriminante de mantenimiento
reutilizado en otra máquina. Cada rechazo se contrastó con su SQLSTATE;
un resultado inesperado detiene la ejecución. La transacción terminó en
\texttt{ROLLBACK}, por lo que los datos sintéticos de prueba no forman
parte del esquema entregado.

También se ejecutaron las cinco pruebas aisladas: clientes y contactos
(cuatro grupos válidos y 167 rechazos), juegos y técnicos (74 rechazos),
premios y canjes (58 casos), sucursales y operaciones (36 resultados),
y personal y administración (once resultados). Las pruebas aisladas usan
dependencias mínimas; la comprobación conjunta utiliza las tablas reales.

Los dominios locales no garantizan reglas operativas entre filas o tablas,
como suficiencia de puntos, descuento de inventario y exclusión de puestos.
Tales reglas requieren transacciones de aplicación u otros mecanismos
adicionales. La validación aquí descrita acredita las restricciones incluidas
en el DDL, sin atribuirles esas reglas operativas.
